\documentclass[11pt,a4paper]{article}

\usepackage[margin=2.5cm]{geometry}
\usepackage{booktabs}
\usepackage{xcolor}
\usepackage{colortbl}
\usepackage{tabularx}
\usepackage{hyperref}
\usepackage{enumitem}
\usepackage{titlesec}
\usepackage{fancyhdr}
\usepackage{array}
\usepackage{microtype}
\usepackage{parskip}
\usepackage{tcolorbox}
\usepackage{fontawesome5}

% ── colours ──────────────────────────────────────────────────────────────────
\definecolor{mainblue}{HTML}{2E5FA3}
\definecolor{lightblue}{HTML}{D6E4F7}
\definecolor{lightgray}{HTML}{F5F5F5}
\definecolor{darkgray}{HTML}{333333}
\definecolor{donegreen}{HTML}{27AE60}
\definecolor{nextorg}{HTML}{E67E22}
\definecolor{planblue}{HTML}{2980B9}
\definecolor{donebg}{HTML}{EAFAF1}
\definecolor{nextbg}{HTML}{FEF9E7}
\definecolor{planbg}{HTML}{EBF5FB}

% ── tcolorbox styles ──────────────────────────────────────────────────────────
\tcbuselibrary{skins, breakable}
\tcbset{
  prioritybox/.style={
    enhanced, breakable,
    colback=#1!8, colframe=#1, arc=3pt,
    boxrule=0.8pt, left=8pt, right=8pt, top=6pt, bottom=6pt,
    fonttitle=\bfseries\color{#1},
  }
}

% ── section formatting ────────────────────────────────────────────────────────
\titleformat{\section}
  {\large\bfseries\color{mainblue}}
  {\thesection.}{0.5em}{}
  [\vspace{-4pt}\textcolor{mainblue}{\rule{\linewidth}{0.8pt}}]

\titleformat{\subsection}
  {\normalsize\bfseries\color{darkgray}}
  {\thesubsection.}{0.5em}{}

\titlespacing{\section}{0pt}{16pt}{6pt}
\titlespacing{\subsection}{0pt}{10pt}{4pt}

% ── header / footer ───────────────────────────────────────────────────────────
\pagestyle{fancy}
\fancyhf{}
\fancyhead[L]{\small\textcolor{mainblue}{\textbf{TCR--Epitope Binding Prediction}}}
\fancyhead[R]{\small\textcolor{darkgray}{Next Steps \& Roadmap}}
\fancyfoot[C]{\small\textcolor{darkgray}{\thepage}}
\renewcommand{\headrulewidth}{0.4pt}

% ── hyperref ──────────────────────────────────────────────────────────────────
\hypersetup{
  colorlinks=true, linkcolor=mainblue, urlcolor=mainblue,
  pdftitle={TCR-Epitope Next Steps},
  pdfauthor={Helena Martinez Rio},
}

% ── status badge macro ────────────────────────────────────────────────────────
\newcommand{\done}{\textcolor{white}{\colorbox{donegreen}{\small\textbf{Done}}}}
\newcommand{\next}{\textcolor{white}{\colorbox{nextorg}{\small\textbf{Next}}}}
\newcommand{\plan}{\textcolor{white}{\colorbox{planblue}{\small\textbf{Planned}}}}

% ─────────────────────────────────────────────────────────────────────────────
\begin{document}

% ── title block ──────────────────────────────────────────────────────────────
\begin{center}
  {\LARGE\bfseries\color{mainblue} TCR--Epitope Binding Prediction}\\[6pt]
  {\large\bfseries\color{darkgray} Next Steps \& Roadmap}\\[8pt]
  {\small\color{gray} VT1 Project \quad|\quad Helena Mart\'inez R\'io \quad|\quad ZHAW 2025}
\end{center}

\vspace{6pt}
\textcolor{mainblue}{\rule{\linewidth}{1.5pt}}
\vspace{12pt}

% ─────────────────────────────────────────────────────────────────────────────
\section{Where We Are}

A classical ML baseline pipeline is complete.
Random Forest with ESM2 embeddings achieves \textbf{AUC 0.738}.
UMAP analysis confirms that the joint TCR+Epitope embedding space has
biologically meaningful structure (pathogen-specific clusters, MHC class separation).
The next phase focuses on improving negative generation, building a proper
evaluation framework, and moving to graph-based models.

% ─────────────────────────────────────────────────────────────────────────────
\section{Roadmap}

\vspace{4pt}
\renewcommand{\arraystretch}{1.4}
\begin{tabularx}{\linewidth}{>{\centering\bfseries\color{mainblue}}p{0.5cm}
                              >{\bfseries}p{3.2cm}
                              X
                              >{\centering\arraybackslash}p{1.5cm}}
\toprule
\rowcolor{mainblue}
\textcolor{white}{\#} &
\textcolor{white}{Step} &
\textcolor{white}{Description} &
\textcolor{white}{Status} \\
\midrule

\rowcolor{donebg}
1 & Improve negatives &
Epitope-aware mismatched pairing: for each positive pair, the negative TCR
comes from a different specificity group, reducing false negative contamination. &
\done \\

\rowcolor{lightgray}
2 & Epitope-split evaluation &
Train/test split via \texttt{GroupShuffleSplit} by epitope so no epitope
appears in both splits. Tests true generalisation to unseen epitopes. &
\done \\

\rowcolor{donebg}
3 & Embedding analysis &
UMAP projections of TCR and epitope embeddings separately.
Silhouette score and intra/inter distance analysis quantify clustering by epitope. &
\done \\

\rowcolor{nextbg}
4 & Larger ESM2 model &
Replace \texttt{esm2\_t6\_8M} with \texttt{esm2\_t12\_35M} or
\texttt{esm2\_t30\_150M} for richer representations.
Expected AUC improvement of 2--5 points. &
\next \\

\rowcolor{lightgray}
5 & Better negatives v2 &
Generate negatives by pairing a TCR with an epitope from a different V-gene
cluster. Biologically motivated and harder for the model to exploit spuriously. &
\next \\

\rowcolor{nextbg}
6 & Knowledge graph integration &
Use the iggytop Neo4j graph to add graph-level scalar features: TCR promiscuity
score, epitope immunodominance score, MHC restriction pathway. &
\next \\

\rowcolor{planbg}
7 & GNN model &
Bipartite graph: TCR nodes + epitope nodes + binding edges.
GraphSAGE or GAT for link prediction.
Compare AUC against RF baseline. &
\plan \\

\rowcolor{lightgray}
8 & Cross-dataset validation &
Evaluate the final model on an external holdout (IEDB or McPAS-TCR)
to test generalisation beyond iggytop. &
\plan \\

\bottomrule
\end{tabularx}

% ─────────────────────────────────────────────────────────────────────────────
\section{Priority for Next Session}

\begin{tcolorbox}[prioritybox=nextorg, title=Step 4 --- Larger ESM2 Model]
The current pipeline uses the smallest ESM2 variant (8M parameters).
Upgrading requires a single config change and a cluster re-run.
This is expected to be the highest-impact single improvement.

\vspace{4pt}
\begin{itemize}[leftmargin=1.2em, itemsep=2pt]
  \item \textbf{Code:} set \texttt{ESM\_MODEL = "facebook/esm2\_t12\_35M\_UR50D"}
        in \texttt{tcr\_epitope\_minimal.py}
  \item \textbf{Cluster:} re-run \texttt{submit\_minimal.slurm} on
        \texttt{sacramento} node (A100 40\,GB)
  \item \textbf{Evaluation:} compare AUC vs current baseline of 0.738
\end{itemize}
\end{tcolorbox}

\vspace{8pt}

\begin{tcolorbox}[prioritybox=planblue, title=Step 6 --- Knowledge Graph Features]
The iggytop database already exposes a Neo4j knowledge graph.
The goal is to extract scalar graph features and add them to the ESM2 feature matrix.

\vspace{4pt}
\begin{itemize}[leftmargin=1.2em, itemsep=2pt]
  \item Run \texttt{create\_knowledge\_graph.py} to populate the local Neo4j instance
  \item Query: for each TCR node, how many distinct epitopes is it connected to?
        (\emph{promiscuity score})
  \item Query: for each epitope node, how many TCRs bind it?
        (\emph{immunodominance score})
  \item Append these scalar features to $X$ alongside ESM2 embeddings and re-train RF
\end{itemize}
\end{tcolorbox}

\vspace{8pt}

\begin{tcolorbox}[prioritybox=mainblue, title=Step 7 --- GNN Architecture Design]
Based on the embedding analysis, the recommended GNN design is:

\vspace{4pt}
\begin{itemize}[leftmargin=1.2em, itemsep=2pt]
  \item \textbf{Graph type:} bipartite --- TCR nodes (features = ESM2 TCR emb.)
        and epitope nodes (features = ESM2 epitope emb.)
  \item \textbf{Edges:} known binding pairs from the training set (positive examples)
  \item \textbf{Task:} link prediction --- does an edge exist between
        TCR$_i$ and epitope$_j$?
  \item \textbf{Architecture:} GraphSAGE (inductive; handles unseen nodes at inference)
        or GAT (attention over neighbours)
  \item \textbf{Library:} PyTorch Geometric (\texttt{torch\_geometric})
\end{itemize}
\end{tcolorbox}

% ─────────────────────────────────────────────────────────────────────────────
\section{Expected Timeline}

\vspace{4pt}
\renewcommand{\arraystretch}{1.3}
\begin{tabularx}{\linewidth}{>{\bfseries\color{mainblue}}p{1.6cm} p{4.5cm} X}
\toprule
\rowcolor{mainblue}
\textcolor{white}{Week} &
\textcolor{white}{Goal} &
\textcolor{white}{Deliverable} \\
\midrule
\rowcolor{lightgray}
Week 1 & Larger ESM2 + re-evaluation &
New AUC benchmark; updated results\_summary.tex \\

Week 2 & Knowledge graph feature extraction &
Graph scalar features added to feature matrix; AUC comparison \\

\rowcolor{lightgray}
Week 3 & GNN prototype (GraphSAGE) &
First GNN AUC vs RF comparison in MLflow \\

Week 4 & GNN tuning + cross-validation &
Final model evaluation; report draft \\
\bottomrule
\end{tabularx}

\vspace{16pt}
\textcolor{mainblue}{\rule{\linewidth}{0.8pt}}
\begin{center}
\small\color{gray}
Generated as part of the VT1 computational immunology project, ZHAW 2025.
\end{center}

\end{document}